\documentclass[11pt]{article}

% ---------- Packages ----------
\usepackage[margin=1in]{geometry}
\usepackage{amsmath,amssymb,amsthm,mathtools}
\usepackage{booktabs}
\usepackage{siunitx}
\usepackage{microtype}
\usepackage{algorithm}
\usepackage{algpseudocode}
\usepackage{hyperref}
\usepackage{graphicx}
\usepackage{xcolor}

% ---------- Metadata ----------
\title{\bf Mantissa Reserve: High-Precision--First Inference for Phase-Faithful Sparse Mixture-of-Experts and Energy-Aware Scheduling}
\author{Payton Ison \and GPT-5 Pro}
\date{November 8, 2025}

% ---------- Macros ----------
\newcommand{\R}{\mathbb{R}}
\newcommand{\E}{\mathbb{E}}
\newcommand{\Var}{\mathrm{Var}}
\newcommand{\Prob}{\mathbb{P}}
\newcommand{\norm}[1]{\left\lVert #1 \right\rVert}
\newcommand{\ip}[2]{\left\langle #1,\,#2 \right\rangle}
\newcommand{\bfx}{\mathbf{x}}
\newcommand{\bfW}{\mathbf{W}}
\newcommand{\bfy}{\mathbf{y}}
\newcommand{\bfe}{\mathbf{e}}
\newcommand{\angles}[2]{\angle\!\left(#1,\,#2\right)}
\newcommand{\ud}{\,\mathrm{d}}

\theoremstyle{plain}
\newtheorem{theorem}{Theorem}
\newtheorem{proposition}{Proposition}
\newtheorem{lemma}{Lemma}

\theoremstyle{definition}
\newtheorem{definition}{Definition}

\begin{document}
\maketitle

\begin{abstract}
We study how floating-point bit allocation---especially mantissa width---governs the stability of \emph{phase information} (directional structure and rankings induced by inner products) in modern inference workloads. Sparse Mixture-of-Experts (MoE) models are an extreme stress test: tiny perturbations to router logits and expert aggregates can flip top-$k$ selection and derail efficiency gains. We formalize \emph{phase fidelity} via angular preservation and rank stability, derive bounds connecting mantissa bits to layerwise phase drift, and show why an ``extra $23$-bit precision layer'' (i.e., FP32 accumulation) qualitatively changes these bounds even when multiplies use FP8/FP16/bfloat16.

We propose \textbf{Mantissa Reserve Accumulation (MRA)}, a drop-in compute pattern that (i) performs low-bit multiplies, (ii) accumulates with a $23$-bit mantissa reserve, and (iii) carries a compact residual to subsequent phase-critical sites. For MoE, we introduce \textbf{Phase-Fidelity Quantization (PFQ)} that concentrates precision at the router and expert-combine boundaries, and we provide sufficient conditions under which top-$k$ remains unchanged. Finally, we present \textbf{PRISM}, a precision-aware datacenter scheduler that allocates the mantissa reserve only where it provably improves phase fidelity under energy/SLO constraints.

Our analysis isolates where precision truly \emph{buys} reliability: not uniformly everywhere, but at a small set of phase-critical reductions. The resulting designs preserve accuracy while enabling aggressive low-bit inference for the majority of FLOPs.
\end{abstract}

\section{Introduction}

The hardware and energy profile of large language model (LLM) inference is increasingly dominated by memory movement and reduction operations, not raw multiply throughput. Precision reduction (FP8/FP16/bfloat16) helps throughput and memory pressure, but it risks corrupting the \emph{phase} of intermediate representations: the direction of vectors, the sign pattern of features, and the ranking of router logits. In sparse Mixture-of-Experts (MoE) models, such phase errors are amplified: a single rank flip can route tokens to the wrong experts, harming both quality and efficiency.

This paper makes three contributions:

\begin{enumerate}
  \item \textbf{A precision--phase theory.} We provide bounds that connect mantissa width to \emph{angular error} and to \emph{rank stability} of inner products. The key constant is unit roundoff $u=2^{-(t+1)}$ for mantissa bits $t$. For dot products of length $n$, the classical $\gamma_n \approx \frac{nu}{1-nu}$ constant governs worst-case relative error; when accumulation uses a $t=23$-bit mantissa (FP32), $\gamma_n$ remains small for practical $n$, while with $t\in\{2,3,7,10\}$ (FP8, bfloat16, FP16) it can exceed $1$ for large $n$, invalidating stability guarantees.
  \item \textbf{Mantissa Reserve Accumulation (MRA).} We propose a compute pattern that retains low-bit multiplies (FP8/FP16/bfloat16) but \emph{always} accumulates with an extra $23$-bit mantissa reserve (FP32). We further store a compact residual (\emph{phase carry}) that is injected exactly at phase-critical barriers (e.g., router logits, attention score reductions, and expert combines).
  \item \textbf{PRISM scheduling for energy efficiency.} We formulate a cluster scheduler that allocates the mantissa reserve where it improves a \emph{Phase Sensitivity Index (PSI)} and satisfies SLOs while minimizing energy. The result is precision \emph{elasticity}: the system is mostly low-bit but ``snaps'' to a high-precision reserve at the few reductions that matter.
\end{enumerate}

We emphasize pragmatism: the method requires only standard FP32 accumulators---ubiquitous in commodity accelerators---and simple software hooks. No specialized hardware is assumed.

\section{Background and Notation}

\paragraph{Floating-point formats.}
We denote a floating-point format by $(s,e,t)$: sign bit $s=1$, exponent width $e$, and mantissa width $t$. For IEEE FP32, $(1,8,23)$; FP16, $(1,5,10)$; bfloat16, $(1,8,7)$; common FP8 variants include E4M3 $(1,4,3)$ and E5M2 $(1,5,2)$. The unit roundoff (maximum relative rounding error to nearest) is $u=2^{-(t+1)}$ for normalized numbers. Thus:
\[
u_{\text{FP32}} = 2^{-24}\approx 5.96\times 10^{-8},\;
u_{\text{FP16}} = 2^{-11} \approx 4.88\times 10^{-4},\;
u_{\text{bf16}} = 2^{-8} = 3.906\times 10^{-3},\;
u_{\text{FP8(E4M3)}} = 2^{-4}=6.25\times 10^{-2}.
\]

\paragraph{Phase fidelity.}
Let $\bfx,\tilde\bfx\in\R^d$ denote a true and a perturbed vector. We quantify phase by:
\begin{itemize}
  \item \emph{Cosine alignment} $A(\bfx,\tilde\bfx)=\frac{\ip{\bfx}{\tilde\bfx}}{\norm{\bfx}\,\norm{\tilde\bfx}}$ and its angle $\angles{\bfx}{\tilde\bfx}=\arccos A(\bfx,\tilde\bfx)$.
  \item \emph{Sign Flip Rate (SFR)} $=\frac{1}{d}\sum_i \mathbf{1}\{\mathrm{sign}(x_i)\neq \mathrm{sign}(\tilde x_i)\}$.
  \item \emph{Rank stability} under inner products: given scores $s_i=\ip{\bfw_i}{\bfx}$, define the top-$k$ set $\mathcal{T}_k$; a perturbation produces $\tilde s_i$ and $\tilde{\mathcal{T}}_k$. The \emph{Gating Stability Index (GSI)} is $\mathrm{GSI}=\Prob[\tilde{\mathcal{T}}_k=\mathcal{T}_k]$ over inputs or noise.
\end{itemize}

\paragraph{Dot-product error constants.}
For a dot product $s=\ip{\bfa}{\bfb}$ of length $n$ under rounding to nearest with unit roundoff $u$, classical analysis yields
\begin{equation}
\label{eq:gamma}
\left|\mathrm{fl}(\ip{\bfa}{\bfb}) - \ip{\bfa}{\bfb}\right|
\;\le\; \gamma_n \sum_{i=1}^n |a_i b_i|,\qquad
\gamma_n := \frac{nu}{1-nu}\,,
\end{equation}
assuming $nu<1$ and no under/overflow during accumulation.
The qualitative behavior is stark: if $nu\ll 1$, the dot product is well conditioned; if $nu\gtrsim 1$, the bound is vacuous.

\section{Why an ``Extra 23 Bits'' Changes the Game}

Even when multiplications are performed at low precision, \emph{accumulation precision} dominates the error in long reductions (e.g., GEMM inner sums, attention score reduces, expert combines). Consider a length-$n$ dot product executed with:
\begin{enumerate}
  \item low-bit multiplies and low-bit accumulation (mantissa $t_\times$, $t_+$), versus
  \item low-bit multiplies with \emph{FP32 accumulation} ($t_+=23$).
\end{enumerate}
Let $u_\times=2^{-(t_\times+1)}$ and $u_+=2^{-(t_++1)}$ be the unit roundoffs for multiply and accumulation, respectively. A coarse yet useful bound is
\[
\left|\mathrm{fl}(\ip{\bfa}{\bfb}) - \ip{\bfa}{\bfb}\right|
\;\lesssim\;
\left(
\underbrace{u_\times}_{\text{product quant.}}
+
\underbrace{\frac{nu_+}{1-nu_+}}_{\text{accumulation}}
\right)\sum_i |a_i b_i|\,.
\]
For $n=4096$ (typical head dimension or MLP width), the terms are illustrative:
\[
\gamma_n^{\text{accum}}(t_+=23)\approx \frac{4096\cdot 2^{-24}}{1-4096\cdot 2^{-24}} \approx 2.44\times 10^{-4}\,,
\]
whereas with $t_+\in\{10,7,3\}$ we get $nu_+\in\{2.0,\;16.0,\;256.0\}$ and the bound becomes useless. This is the essence of the \emph{mantissa reserve}: a $23$-bit accumulator qualitatively restores a small $\gamma_n$ for realistic $n$ even if the inputs are low-bit.

\section{From Magnitude to Phase: Angular and Rank Stability}

\subsection{Angular error across linear layers}

Let $\tilde\bfx=\bfx+\bfe$ with $\norm{\bfe}\le \varepsilon\norm{\bfx}$. Then
\begin{equation}
\label{eq:angle-bound}
1 - A(\bfx,\tilde\bfx) \;\le\; \tfrac{1}{2}\varepsilon^2 + O(\varepsilon^3),\qquad
\angles{\bfx}{\tilde\bfx} \;\le\; \varepsilon + O(\varepsilon^2)\,.
\end{equation}
For a linear map $\bfy=\bfW\bfx$, $\tilde\bfy=\mathrm{fl}(\bfW)\tilde\bfx+\delta$ where $\delta$ collects dot-product errors from accumulation. If each output coordinate is a length-$n$ reduction with accumulation unit roundoff $u_+$, then, neglecting second-order terms,
\[
\norm{\delta}_\infty \lesssim \gamma_n^{\text{accum}} \max_j \sum_i |W_{ji} x_i|\,.
\]
When accumulation uses the mantissa reserve ($t_+=23$), $\gamma_n$ is small, and the resulting \emph{angular drift} bound follows by a standard perturbation argument:
\begin{proposition}[Layerwise phase drift]
\label{prop:phase-drift}
Ignoring under/overflow and assuming normalized activations $\norm{\bfx}=1$, the expected cosine misalignment after a linear layer satisfies
\[
1 - \E\!\left[A(\bfy,\tilde\bfy)\right]
\;\lesssim\;
c_1 \, u_\times^2 + c_2 \, \left(\gamma_n^{\text{accum}}(t_+{=}23)\right)^2,
\]
for constants $c_1,c_2$ depending on weight statistics. The dominant term is the accumulation constant whenever $n$ is large.
\end{proposition}

\subsection{Rank stability for gating (MoE routers)}

Let $s_i=\ip{\bfw_i}{\bfx}$ denote router logits, and let $m_k:=s_{(k)}-s_{(k+1)}$ be the top-$k$ margin (order statistics). Low-bit multiplies plus high-precision accumulation yield $\tilde s_i=s_i+\Delta_i$ with $|\Delta_i|\le \eta:=\gamma_n^{\text{accum}} \max_i \norm{\bfw_i}\norm{\bfx} + u_\times \sum_j |w_{ij} x_j|$. Then:

\begin{theorem}[Sufficient condition for top-$k$ stability]
\label{thm:topk}
If $m_k > 2\eta$, then the top-$k$ index set is unchanged: $\tilde{\mathcal{T}}_k=\mathcal{T}_k$.
\end{theorem}

\noindent
\textbf{Implication.} For realistic $n$, the $\eta$ term is dominated by the accumulation constant; using a $23$-bit accumulator (even with FP8/FP16 multiplies) sharply increases the regime where $m_k>2\eta$, explaining empirical observations that ``FP8 with FP32 accumulate'' preserves router choices while ``FP8 throughout'' does not.

\section{Mantissa Reserve Accumulation (MRA)}
\label{sec:mra}

\subsection{Design goals}
\begin{itemize}
  \item Retain low-bit multiplies to save bandwidth and memory.
  \item Use FP32 accumulation wherever phase is at risk: long reductions, router logits, attention score reductions, expert combines, layernorm statistics.
  \item Track and re-inject a compact \emph{phase carry} residual around non-linearities to reduce drift across blocks.
\end{itemize}

\subsection{Computation pattern}

Let $A\in\R^{m\times n}$, $B\in\R^{n\times p}$ be activations/weights stored at low precision (e.g., FP8 or bfloat16). MRA performs:
\begin{enumerate}
  \item Dequantize tiles to a fused multiply representation at low precision (\texttt{FP8/FP16}).
  \item Multiply at low precision but accumulate into FP32 registers (FMA rounds once).
  \item Apply a compensated reduction (Kahan/Neumaier-style) to further suppress cancellation when $n$ is large.
  \item Optionally cache a \emph{phase carry} $r$ (the compensation term) and re-inject $r$ at subsequent phase-critical boundaries.
\end{enumerate}

\begin{algorithm}[t]
\caption{MRA: GEMM micro-kernel (schematic)}
\label{alg:mra}
\begin{algorithmic}[1]
\Require Tiles $A,B$ in low-bit format; output tile $C$ in FP32
\State $C \gets 0$ \Comment{FP32 accumulator tile}
\State $R \gets 0$ \Comment{phase-carry (compensation)}
\For{each $k$-tile}
  \State $\hat A \gets \mathrm{dequantize}(A_{\cdot k})$ \Comment{kept in low-bit lanes}
  \State $\hat B \gets \mathrm{dequantize}(B_{k \cdot})$
  \State $Y \gets \hat A \otimes \hat B$ \Comment{low-bit multiply}
  \State \textbf{[FP32 accumulate with compensation]}
  \State $y \gets \mathrm{sum}(Y)$ \Comment{tile reduce in FP32}
  \State $t \gets C + y$
  \State $R \gets (t - C) - y + R$
  \State $C \gets t$
\EndFor
\State \Return $C + R$ \Comment{inject phase carry}
\end{algorithmic}
\end{algorithm}

\paragraph{Why carry a residual?}
Many phase-critical barriers are followed by non-linearities (gating softmax, layernorm, activation functions) that magnify small directional errors. Passing a small FP32 residual to the next barrier materially reduces cumulative angular drift over a block without paying FP32 throughout.

\section{Phase-Fidelity Quantization (PFQ) for MoE}
\label{sec:pfq}

MoE introduces two fragile reductions:
\begin{enumerate}
  \item \textbf{Routing.} Top-$k$ selection on router logits; requires rank stability (\S\ref{thm:topk}).
  \item \textbf{Expert combine.} Weighted sum of expert outputs; requires angular stability of the aggregate.
\end{enumerate}
PFQ concentrates mantissa reserve at these edges while leaving per-expert MLPs largely low-bit internally.

\subsection{PFQ recipe}
\begin{itemize}
  \item \emph{Router logits:} compute $\ip{\bfw_i}{\bfx}$ with FP32 accumulation; store logits in FP16/bfloat16 only after top-$k$ is decided.
  \item \emph{Dispatch weights:} compute softmax in FP32 (or \textit{log-sum-exp} in FP32) to preserve small gaps $m_k$; downcast weights post-normalization.
  \item \emph{Expert outputs:} experts may compute in FP8/FP16 with FP32 accumulate \emph{inside} each MLP; the final expert \emph{combine} uses FP32 accumulation across experts.
  \item \emph{Normalization/attention:} statistics (mean/var; $\mathrm{softmax}$ denominator) use FP32 accumulation.
\end{itemize}

\subsection{Sufficient conditions for robust routing}

Combining Theorem~\ref{thm:topk} with a margin model (e.g., sub-Gaussian router scores) yields a conservative bound on $\mathrm{GSI}$:
\[
\mathrm{GSI} \;\ge\; 1 - \Prob\!\left[m_k \le 2\eta\right]
\quad\text{with}\quad
\eta \approx \gamma_n^{\text{accum}}(23)\cdot \Lambda + u_\times \cdot \Lambda',
\]
where $\Lambda,\Lambda'$ capture scale factors (norms of inputs/weights). In practice, the $23$-bit accumulator renders $\gamma_n$ the small term even for large $n$, boosting $\mathrm{GSI}$.

\section{PRISM: Precision-Aware Scheduling for Energy/SLO}
\label{sec:prism}

\subsection{Problem setup}
Consider a set of inference requests $\mathcal{J}$ with deadlines $d_j$ and per-request phase sensitivity $\mathrm{PSI}_j$ (estimated from router margins, attention entropy, or calibration statistics). Let $\mathcal{P}$ be precision profiles (e.g., \{FP8$\times$ / FP32$+$, bf16$\times$ / FP32$+$, FP16$\times$ / FP32$+$\}) and $\mathcal{M}$ be machine types. Each tuple $(m,p)$ has energy-per-token $E_{m,p}$ and service time $T_{m,p}$.

\subsection{Objective}
Assign $(m,p)$ and start times to minimize total energy while meeting SLOs and ensuring that phase-critical segments run with the mantissa reserve when $\mathrm{PSI}$ exceeds a threshold.

\paragraph{Mixed-integer formulation (sketch).}
Let $x_{j,m,p}\in\{0,1\}$ indicate assignment and $\tau_j$ start time. The objective is
\[
\min \sum_{j}\sum_{m,p} E_{m,p}\, x_{j,m,p}
\]
subject to capacity constraints, $\,\tau_j + T_{m,p}\le d_j$ if $x_{j,m,p}=1$, and a \emph{phase constraint}
\[
\mathrm{PSI}_j > \theta \;\Rightarrow\; p \in \mathcal{P}_{\text{reserve}},
\]
where $\mathcal{P}_{\text{reserve}}$ includes profiles with FP32 accumulation at phase-critical barriers.

\subsection{A practical heuristic}
PRISM uses a two-stage greedy allocator:
\begin{enumerate}
  \item \textbf{Phase-first filtering.} If $\mathrm{PSI}_j>\theta$, require a reserve-capable profile (keeps the router/combines in FP32 accumulate).
  \item \textbf{Energy-aware bin packing.} Sort remaining jobs by laxity $(d_j - \hat T_{j})$ and assign the lowest-energy profile that meets SLOs; co-schedule low-$\mathrm{PSI}$ jobs on aggressively low-bit profiles.
\end{enumerate}
This achieves most of the energy benefit while honoring the small set of phase-critical cases.

\section{Practical Guidance}

\paragraph{Where to spend precision.}
\begin{enumerate}
  \item Accumulators in long reductions (GEMM inner sums).
  \item Router logits and expert combines in MoE.
  \item Attention softmax denominators and layernorm statistics.
\end{enumerate}

\paragraph{Where to save.}
\begin{itemize}
  \item Expert-internal multiplies: FP8/FP16/bfloat16 multiplies are typically safe if accumulation is FP32.
  \item Activation storage between non-critical layers.
\end{itemize}

\paragraph{Numerical hygiene.}
Use fused operations (FMA), avoid subtractive cancellation (e.g., prefer \textit{log-sum-exp}), and enable compensated summation on exceptionally long reductions.

\section{Illustrative Calculations}
\label{sec:calc}

Consider $n=4096$ and compare accumulation precisions (keeping multiplies low-bit):
\begin{center}
\begin{tabular}{lccc}
\toprule
Accumulator & mantissa $t_+$ & $u_+$ & $\gamma_n \approx \frac{nu_+}{1-nu_+}$ \\
\midrule
FP32 & 23 & $5.96\!\times\!10^{-8}$ & $2.44\!\times\!10^{-4}$ \\
FP16 & 10 & $4.88\!\times\!10^{-4}$ & $>1$ (no guarantee) \\
bfloat16 & 7 & $3.91\!\times\!10^{-3}$ & $>1$ (no guarantee) \\
FP8 (E4M3) & 3 & $6.25\!\times\!10^{-2}$ & $>1$ (no guarantee) \\
\bottomrule
\end{tabular}
\end{center}

Even with FP8 multiplies, the FP32 accumulator keeps $\gamma_n$ small enough to preserve router rankings under moderate margins (Theorem~\ref{thm:topk}). Without it, worst-case guarantees fail for typical widths.

\section{Threats to Validity and Limitations}
Our bounds are worst-case and conservative; real workloads enjoy structure (sparsity, normalization) that is not fully exploited here. We do not account for denormal behavior, subnormal flushing, or stochastic rounding. Hardware-specific behaviors (e.g., tensor core accumulation paths) must be validated per platform. Finally, \emph{phase carry} requires minimal but nonzero extra state; the best compression strategy is workload-dependent.

\section{Related Work (brief)}
Mixed-precision training/inference leverages FP16/bfloat16 and FP32 accumulation widely. FP8 formats have been explored for both training and inference; it is now standard to accumulate sensitive reductions in FP32. Sparse MoE architectures (e.g., sparsely-gated and Switch-style routers) underscore the importance of router stability. Our contribution is to formalize \emph{phase fidelity}, tie it explicitly to mantissa-driven accumulation constants, and systematize where a $23$-bit reserve is \emph{necessary and sufficient} for stability, along with a scheduler that deploys it only where it pays.

\section{Conclusion}

Precision is a \emph{budget} to be spent at a few critical reductions rather than a uniform tax on all operations. The ``extra $23$ bits'' of an FP32 accumulator collapse the accumulation constant $\gamma_n$ and thereby stabilize the directions and rankings that govern routing and global reductions. \textbf{MRA} and \textbf{PFQ} convert this observation into implementable kernels and model recipes, while \textbf{PRISM} allocates the mantissa reserve under energy/SLO constraints. Together, these tools preserve accuracy for sparse MoE inference while enabling aggressive low-bit execution almost everywhere else.

\paragraph{Reproducibility.} The algorithms above require only standard BLAS-like interfaces with FP32 accumulators and can be integrated via micro-kernel updates. We recommend logging router margins and attention entropies to estimate $\mathrm{PSI}$ for PRISM.

\section*{Acknowledgments}
We thank colleagues for helpful discussions on numerical stability and systems design. Any remaining errors are our own.

% ----------------- Appendix -----------------
\appendix

\section{Appendix A: Proof Sketch of Theorem~\ref{thm:topk}}

Let $s_{(1)}\ge \cdots \ge s_{(k)} \ge s_{(k+1)} \ge \cdots$ be ordered scores and suppose $|\Delta_i|\le \eta$ for all $i$. For any $i\in \mathcal{T}_k$ and $j\notin \mathcal{T}_k$, we have $s_i \ge s_{(k)}$ and $s_j \le s_{(k+1)}$, so
\[
\tilde s_i - \tilde s_j \;\ge\; (s_i - s_j) - 2\eta \;\ge\; m_k - 2\eta\,.
\]
If $m_k > 2\eta$, then $\tilde s_i > \tilde s_j$, which preserves all cross-boundary pairwise orderings. Hence $\tilde{\mathcal{T}}_k=\mathcal{T}_k$.

\section{Appendix B: Compensated Summation for Long Reductions}

Kahan summation maintains $(\mathrm{sum},\mathrm{comp})$ such that the rounding error from each addition is tracked and re-injected. In MRA we implement a vectorized variant, which empirically reduces phase drift in deep reductions (notably attention and expert combine) without material overhead because the compensation state fits in FP32 registers.

\section{Appendix C: Phase Sensitivity Index (PSI)}

Define per-token $\mathrm{PSI}$ as a convex combination of normalized indicators:
\[
\mathrm{PSI} \;=\; \alpha \cdot \frac{1}{\log k}\,H(\mathrm{softmax}(s))
\;+\; \beta \cdot \frac{\norm{\nabla_{\bfx}\, \mathrm{logit}}_2}{\mathrm{EMA}}
\;+\; \gamma \cdot \mathrm{SNR}^{-1}\!,
\]
where $H$ is entropy of router logits, the gradient norm proxies sensitivity, and $\mathrm{SNR}$ comes from calibration statistics. Higher $\mathrm{PSI}$ implies more precision benefit at phase-critical sites.

\section*{References}
\vspace{-0.5em}
\begin{itemize}\setlength\itemsep{0.25em}
  \item IEEE Standard for Floating-Point Arithmetic (IEEE 754-2019).
  \item N. J. Higham. \textit{Accuracy and Stability of Numerical Algorithms}, 2nd ed., SIAM, 2002.
  \item S. Shazeer et al. ``Outrageously Large Neural Networks: The Sparsely-Gated Mixture-of-Experts Layer,'' 2017.
  \item W. Fedus, B. Zoph, N. Shazeer. ``Switch Transformers: Scaling to Trillion Parameter Models with Simple and Efficient Sparsity,'' 2021.
  \item P. Micikevicius et al. ``Mixed Precision Training,'' 2017.
  \item N. P. Jouppi et al. ``In-Datacenter Performance Analysis of a Tensor Processing Unit,'' 2017.
\end{itemize}

\end{document}

